\documentclass[openany,fontset=fandol]{ctexbook}
\usepackage{amsmath,amssymb,tikz-cd}
\usepackage[unicode,hidelinks]{hyperref}
\setlength{\emergencystretch}{2em}
\newcommand{\Oa}{\mathcal O}
\begin{document}
\title{代数几何基础第五卷预备笔记\\可逆模截面的类比图}
\author{据英译编辑稿翻译，参校法文扫描}
\date{2026年11月20日}
\maketitle
\chapter*{范围与接续}
本批接英译第68页引图句，完整译出式（D）的类比图及随后的英译者方括号说明。对读法文物理第103--104页。下一起点是英语同页“Upon arriving at this point”及法文第104页“Arrivé à ce point”。本稿独立增补，未并入第十七版。
\chapter*{可逆模截面的类比图（续）}
下面给出一个与（D）类似的图，不过这里取可逆模$K$的一个截面$\varepsilon$。
\[
\begin{tikzcd}[column sep=0.8em,row sep=1.8em]
 & & 0 & & \\
 & & \Omega^1_{X/Y}\otimes\Oa_{V'} \arrow[u] & & \\
0 \arrow[r] &
 (\mathcal P^1_{X/Y}(K))^\vee\otimes\Oa_{V'}
 \arrow[r] \arrow[ur,"\mu"] &
 \Omega^1_{X/S}\otimes\Oa_{V'} \arrow[u] \arrow[r] &
 \Omega^1_{V'/S} \arrow[r] & 0 \\
 & K^{-1}\otimes\Oa_{V'} \arrow[u] \arrow[r] &
 \Omega^1_{Y/S}\otimes\Oa_{V'} \arrow[u] \arrow[ur,"\nu"'] & & \\
 & \omega_{X/Y}^{-2}\otimes K^{-1}\otimes\Oa_{V'} \arrow[u] &
 0 \arrow[u] & &
\end{tikzcd}
\]

\noindent [以下部分被格罗滕迪克划去，英译者不确定原因。为求完整，仍将其转载；到划去部分结束之处，将标明“不再是划去内容”。]
\chapter*{译校说明（非原文）}
\paragraph{图的对读。}
法文第103页引图的手写截面符号看起来作$\varphi$；本批重绘沿用英译的$\varepsilon$，不因此另设两个截面。图中$v$为对偶符号，重绘用$\vee$。$\mu$是先由左中项入环境微分模、再映至相对微分模的复合；$\nu$由下中项经环境微分模映至右端。图中的顶端和底端零对象属于环境模的竖直正合列，并非另为斜箭头断言单射或满射。

\paragraph{左下角的指数。}
重绘照英译保留$K^{-1}$；法文第103页对应手写指数有涂抹，不能据此确认完整字样。下面只作类型核对。若仍沿用前文
$L=\det P\otimes\det Q\otimes(\det E)^{-1}$，且$P$为秩$d+1$的相对一阶主部模$\mathcal P^1_{X/Y}(K)$之对偶，那么在主部模正合列
\[
0\to\Omega^1_{X/Y}\otimes K
\to\mathcal P^1_{X/Y}(K)\to K\to0
\]
的约定下，
$\det P\simeq \omega_{X/Y}^{-1}\otimes K^{-(d+1)}$。
取$Q=\Omega^1_{Y/S}|_{V'}$和$E=\Omega^1_{X/S}|_{V'}$，将这些线丛拉回并限制至$V'$，传递性行列式关系遂给出
\[
L\simeq
(\omega_{X/Y}^{-2}\otimes K^{-(d+1)})|_{V'}.
\]
因此，若左下角意在表示这条$L$，则指向$K^{-1}|_{V'}$的箭头相应属于
$\operatorname{Hom}(L,K^{-1}|_{V'})
\simeq(\omega_{X/Y}^{2}\otimes K^{d})|_{V'}$，
其$K$幂次与上一批单位变换的$u^d$相容。这不是对模糊手写指数的辨读，也不是在译文中静默更换英译图的标签。

\paragraph{划去标记。}
方括号是英译者说明，不是作者正文。法文第104页从下一段已可见贯页划线；本批只到说明为止，未查定整段划去稿的最终边界。
\end{document}
